% =============================================================================
% 4. fejezet — Mérnöki rész: a DermAI alkalmazás
% =============================================================================

Ez a fejezet a DermAI rendszer mérnöki megvalósítását mutatja be: a betanított
osztályozó modellt egy teljes, végponttól végpontig működő alkalmazássá
egészíti ki, amely okostelefonon és webböngészőben egyaránt elérhető. Először a
követelményeket és a rendszerarchitektúrát ismertetjük, majd komponensenként a
megvalósítást: a modellt kiszolgáló réteget, a backendet, a mobil- és
webalkalmazást, a hitelesítést és a felhőalapú adattárolást, végül a hálózati
hibatűrést és a biztonsági megfontolásokat.

% =============================================================================
\section{Követelmények}
% =============================================================================

\subsection{Funkcionális követelmények}

A rendszerrel szemben támasztott fő funkcionális elvárások a következők:

\begin{itemize}
  \item A felhasználó fényképet készíthet vagy tölthet fel egy
  bőrelváltozásról (kamera vagy galéria).
  \item A rendszer a feldolgozás előtt ellenőrzi a kép minőségét és
  relevanciáját, és elutasítja az értékelhetetlen felvételeket.
  \item A bőrelváltozást egy ViT-alapú modell osztályozza
  (jó\-indulatú / rosszindulatú), konfidencia-értékkel.
  \item A rendszer természetes nyelvű magyarázatot ad az eredményről, az
  ABCDE-szabály figyelembevételével.
  \item A felhasználó utólagos kérdéseket tehet fel az elváltozással
  kapcsolatban (párbeszéd).
  \item A felhasználó bejelentkezik, és a feltöltött képek, valamint a
  beszélgetések a saját fiókjához kötve, biztonságosan tárolódnak.
  \item A felhasználó visszanézheti korábbi munkameneteit (előzmények).
\end{itemize}

\subsection{Nem-funkcionális követelmények}

\begin{itemize}
  \item \textbf{Reszponzivitás:} az AI-válaszok token-szintű, fokozatos
  megjelenítése, hogy a felhasználó ne egy hosszú várakozás után, hanem
  folyamatosan kapjon visszajelzést.
  \item \textbf{Hibatűrés:} a rendszer alkalmazkodjon a változó hálózati
  körülményekhez és a modellt kiszolgáló komponens időszakos
  elérhetetlenségéhez.
  \item \textbf{Biztonság:} a felhasználói adatok hitelesített hozzáférés mögött,
  felhasználónként elkülönítve tárolódjanak; a titkos kulcsok ne kerüljenek a
  forráskódba.
  \item \textbf{Hordozhatóság:} a modellt kiszolgáló komponens címe futás közben
  módosítható legyen, hogy a (gyakran változó címen elérhető) GPU-háttér
  újraindítása ne igényelje az alkalmazás újratelepítését.
\end{itemize}

% =============================================================================
\section{Rendszerarchitektúra}
% =============================================================================

A DermAI több, jól elkülönülő komponensből álló, elosztott rendszer. A két
kliens (Android alkalmazás és webes felület) külön-külön backendhez kapcsolódik,
amelyek azonban közös elveken működnek, és ugyanazt a modellt kiszolgáló réteget,
illetve felhőszolgáltatásokat használják. A komponensek és kapcsolataik
áttekintését a \ref{fig:architecture}.~ábra mutatja.

\begin{itemize}
  \item \textbf{Android alkalmazás} (Kotlin) -- a mobil kliens; a
  \textbf{Flask} alapú backendhez \cite{flask} kapcsolódik.
  \item \textbf{Webes felület} (HTML/JavaScript) -- a böngészőkliens; a
  \textbf{FastAPI} alapú backendhez \cite{fastapi} kapcsolódik.
  \item \textbf{Modellt kiszolgáló réteg} -- egy MCP-szerver
  \cite{mcp2024}, amely a GPU-n futó ViT modellt szólítja meg.
  \item \textbf{Gemini 2.5 Flash} \cite{gemini2023} -- a képvalidációt, a
  magyarázatgenerálást és a párbeszédet végző nagy nyelvi modell.
  \item \textbf{Firebase} \cite{firebase} -- hitelesítés (Authentication),
  adatbázis (Firestore) és fájltárolás (Storage).
\end{itemize}

% MEGJEGYZÉS (Ervin): ide készíts egy komponens-/architektúradiagramot
% (Android + Web -> Flask/FastAPI -> MCP -> Colab ViT; valamint Firebase és Gemini).
% Pl. draw.io-val rajzolt ábra illeszthető ide.
\begin{figure}[H]
  \centering
  \fbox{\begin{minipage}[c][5cm][c]{0.9\textwidth}
    \centering\itshape
    [Ábra helye] -- A DermAI rendszer komponensdiagramja: a két kliens
    (Android, web), a hozzájuk tartozó backendek (Flask, FastAPI), a modellt
    kiszolgáló MCP-réteg és a GPU-n futó ViT, valamint a Firebase és a Gemini
    szolgáltatások kapcsolatai. (Saját szerkesztésű ábrát illessz be.)
  \end{minipage}}
  \caption{A DermAI rendszer magas szintű architektúrája.}
  \label{fig:architecture}
\end{figure}

\subsection{Egy elemzési kérés életciklusa}

Egy kép elemzésének tipikus folyamata a következő lépésekből áll:

\begin{enumerate}
  \item A kliens elküldi a képet és a hitelesítő tokent a backendnek.
  \item A backend ellenőrzi a tokent, majd a képet a Gemini modellnek továbbítja
  \textit{minőség- és relevanciaellenőrzésre}. Ha a kép nem értékelhető (pl.
  nem bőrelváltozás, túl távoli, életlen vagy több léziót tartalmaz), a folyamat
  itt megszakad, és a felhasználó célzott visszajelzést kap.
  \item Megfelelő kép esetén a backend feltölti a képet a felhasználó
  tárhelyére (Firebase Storage), majd az MCP-rétegen keresztül lefuttatja a
  ViT-osztályozást.
  \item A backend létrehozza vagy frissíti a munkamenet rekordját a Firestore
  adatbázisban.
  \item A backend a klasszifikációs eredményt belefoglalja a Gemini számára
  összeállított utasításba, és a magyarázatot \textbf{token-szintű
  adatfolyamként} (SSE) küldi vissza a kliensnek.
\end{enumerate}

% =============================================================================
\section{A modellt kiszolgáló komponens: MCP-szerver}
% =============================================================================

A betanított ViT modell a backendtől elkülönített, GPU-val rendelkező
környezetben (Google Colab) fut, és HTTP-végponton fogadja a
klasszifikációs kéréseket. A backend és e modellkiszolgáló közötti
kommunikációt egy MCP- (Model Context Protocol) szerver \cite{mcp2024} közvetíti,
amely a backendből indított alfolyamatként, szabványos be- és kimeneti (stdio)
csatornán, JSON-RPC 2.0 üzenetekkel \cite{jsonrpc2013} kommunikál. Egy
klasszifikációs kérés üzenetformátuma például:

\begin{verbatim}
{"jsonrpc": "2.0", "id": 1, "method": "tools/call",
 "params": {"name": "classify_lesion",
            "arguments": {"image_base64": "..."}}}
\end{verbatim}

Az MCP-szerver a kapott, base64 kódolású képet továbbítja a GPU-n futó modell
\texttt{/classify} végpontjának, és a választ -- a címkét, a konfidenciát és az
osztályonkénti pontszámokat -- adja vissza a backendnek. A megoldás egyik
kulcseleme egy \texttt{update\_model\_url} metódus, amely lehetővé teszi a
modellkiszolgáló címének \textbf{futás közbeni} módosítását: mivel a GPU-háttér
újraindításakor a nyilvános cím (alagút) megváltozik, e mechanizmus révén a
backend új cím esetén sem igényel újraindítást.

% =============================================================================
\section{A backend}
% =============================================================================

A rendszer két backendet tartalmaz, amelyek a két klienst szolgálják ki, de
azonos belső felépítést követnek. A mobil kliens egy \textbf{Flask}
\cite{flask} alapú szervert, a webes felület pedig egy \textbf{FastAPI}
\cite{fastapi} alapú szervert használ. Mindkét backend ugyanazt az MCP-klienst,
ugyanazt a Gemini-integrációt és ugyanazokat a Firebase-szolgáltatásokat
használja.

\subsection{Végpontok}

A webes (FastAPI) backend főbb végpontjait a \ref{tab:endpoints}.~táblázat
foglalja össze; a mobil (Flask) backend ezekkel funkcionálisan egyenértékű
végpontokat biztosít.

\begin{table}[H]
  \centering
  \caption{A webes backend (FastAPI) főbb végpontjai.}
  \label{tab:endpoints}
  \begin{tabular}{lll}
    \hline
    \textbf{Metódus} & \textbf{Útvonal} & \textbf{Funkció} \\
    \hline
    POST & \texttt{/analyze}            & Kép elemzése (validálás, osztályozás, magyarázat) \\
    POST & \texttt{/chat}               & Utólagos kérdés (kép nélkül), párbeszéd \\
    GET  & \texttt{/history/\{id\}}     & Egy korábbi munkamenet lekérése \\
    GET  & \texttt{/sessions}           & A felhasználó munkameneteinek listája \\
    PUT  & \texttt{/model-url}          & A modellkiszolgáló címének frissítése \\
    GET  & \texttt{/health}             & Állapotellenőrzés \\
    \hline
  \end{tabular}
\end{table}

\subsection{A Gemini integrációja és a képvalidáció}

A nagy nyelvi modell (Gemini 2.5 Flash) a rendszerben három feladatot lát el. Az
\textbf{elsődleges szűrés} során a feltöltött képet egy rögzített utasítással
értékeli, amely egyetlen kódot kér vissza: \texttt{OK} (értékelhető),
\texttt{NOT\_SKIN}, \texttt{TOO\_FAR}, \texttt{MULTIPLE}, \texttt{BLURRY} vagy
\texttt{NOT\_ASSESSABLE}. Ha a kód nem \texttt{OK}, a felhasználó a hibának
megfelelő, célzott útmutatást kap (például: tartsa közelebb a kamerát, vagy
fényképezzen egyetlen léziót). A \textbf{magyarázatgenerálás} során a Gemini a
ViT klasszifikációs eredményét beépítve, az ABCDE-szabály szerint értelmezi az
elváltozást, közérthető nyelven. Végül a Gemini kezeli a felhasználó
\textbf{utólagos kérdéseit} is, a beszélgetés előzményeinek figyelembevételével.

\subsection{Token-szintű adatfolyam (SSE)}

A reszponzív felhasználói élmény érdekében a Gemini válaszait nem egyben, hanem
token-szintű adatfolyamként továbbítjuk a kliensnek, szerver-küldött események
(Server-Sent Events) formájában \cite{whatwg_sse}. Az adatfolyam egyes
eseményei egy-egy szövegdarabot tartalmaznak, a lezárást pedig egy különleges
jelző jelzi:

\begin{verbatim}
data: {"token": "Az elváltozás "}
data: {"token": "aszimmetrikus, ..."}
data: [DONE]
\end{verbatim}

Így a felhasználó a szöveget a generálással egyidejűleg, folyamatosan látja
megjelenni.

% =============================================================================
\section{A mobilalkalmazás}
% =============================================================================

A mobil kliens natív Android alkalmazás, Kotlin nyelven. Indításkor egy
bejelentkező képernyő (\texttt{LoginActivity}) fogadja a felhasználót, ahol
Google-fiókkal jelentkezhet be; sikeres hitelesítés után a fő képernyőre
(\texttt{MainActivity}) jut.

A fő képernyő chat-alapú felületet biztosít: a felhasználó a kamerával
fényképet készíthet (egy dedikált \texttt{CameraActivity} segítségével), vagy a
galériából választhat képet, majd szöveges üzenettel együtt elküldheti azt
elemzésre. A beérkező AI-válaszokat a felület token-szintű adatfolyamként,
folyamatosan jeleníti meg, a klasszifikáció eredményét pedig egy színkódolt
jelvény (jó\-indulatú / rosszindulatú) emeli ki. A menüből elérhető a korábbi
munkamenetek előzménye, az alkalmazás ismertetője (About), a hálózati
beállítások, valamint a kijelentkezés.

\subsection{Hálózati réteg}

Az alkalmazás hálózati rétege az OkHttp és a Retrofit könyvtárakra épül. A
streaming (SSE) válaszokat egy Kotlin \texttt{Flow} kezeli, az állapotot pedig
egy \texttt{ChatViewModel} tartja karban, így az adatfolyam a képernyő-forgatást
is túléli. A backend címét egy beállítási rétegben (\texttt{NetworkConfig})
tároljuk, amely a \texttt{SharedPreferences}-ben őrzi a felhasználó által
megadott vagy az alapértelmezett címet -- ez teszi lehetővé, hogy az alkalmazás
mindig ott keresse a backendet, ahol az ténylegesen elérhető.

% =============================================================================
\section{A webes felület}
% =============================================================================

A webes felület egy egyoldalas, sötét témájú alkalmazás, amely a FastAPI backend
által kiszolgált sablonra épül. A felület a Firebase JavaScript SDK-t használja a
Google-alapú bejelentkezéshez: amíg a felhasználó nincs bejelentkezve, egy
bejelentkező réteg takarja a felületet, sikeres hitelesítés után pedig
megjelenik a chat-felület. Minden, a backendhez intézett kérés tartalmazza a
felhasználó hitelesítő tokenjét. A felületről elérhető egy külön
ismertető (About) oldal, valamint a kijelentkezés is. A token-szintű
válaszokat a böngésző a szerver-küldött események feldolgozásával jeleníti meg
fokozatosan.

% =============================================================================
\section{Hitelesítés és felhőalapú adattárolás}
% =============================================================================

A rendszer a Firebase szolgáltatásait \cite{firebase} használja a
hitelesítéshez és az adattároláshoz.

\subsection{Hitelesítés}

A bejelentkezés mindkét platformon Google-fiókkal, a Firebase Authentication
segítségével történik. A kliens a sikeres bejelentkezés után egy
azonosító\-tokent (ID token) kap, amelyet minden backend-kéréshez mellékel az
\texttt{Authorization} fejlécben. A backend e tokent szerveroldalon ellenőrzi a
Firebase Admin SDK segítségével, és a kérésből kinyeri a felhasználó egyedi
azonosítóját (\texttt{uid}); érvénytelen vagy hiányzó token esetén a kérést
elutasítja.

\subsection{Adatbázis és fájltárolás}

A munkameneteket a Firestore adatbázis tárolja, munkamenetenként egy
dokumentumban, amely tartalmazza a tulajdonos azonosítóját, a klasszifikáció
eredményét, a feltöltött kép azonosítóját és a beszélgetés üzeneteit. A
dokumentum szerkezete vázlatosan:

\begin{verbatim}
sessions/{session_id}
   user_id        : a tulajdonos Firebase-azonosítója
   created_at     : létrehozás időbélyege
   source         : "web" vagy "android"
   image_id       : a feltöltött kép azonosítója
   classification : {label, confidence, all_scores}
   messages       : [{role, text, timestamp}, ...]
\end{verbatim}

A feltöltött képeket a Firebase Storage tárolja, \textbf{felhasználónként
elkülönített} útvonalon: \texttt{users/\{uid\}/lesions/\{image\_id\}.jpg}. Így a
képek a feltöltő felhasználó fiókjához kötődnek, és a munkamenetek listázásakor a
backend a tulajdonos azonosítója szerint szűr, biztosítva, hogy minden
felhasználó kizárólag a saját adatait érje el.

% =============================================================================
\section{Hálózati robusztusság és hibatűrés}
% =============================================================================

Mivel a rendszer több, egymástól független szolgáltatásra támaszkodik -- és a
modellt kiszolgáló GPU-háttér címe is változhat --, a hálózati hibatűrés
kiemelt tervezési szempont volt. A megvalósított mechanizmusok:

\begin{itemize}
  \item \textbf{Újrapróbálkozás exponenciális visszalépéssel.} A REST-hívásokhoz
  egy OkHttp-elfogó (interceptor), az SSE-adatfolyamhoz pedig a Kotlin
  \texttt{Flow} \texttt{retryWhen} operátora biztosít automatikus
  újrapróbálkozást átmeneti hálózati hiba esetén.
  \item \textbf{Az MCP-kliens önjavítása.} A backend MCP-kliense a sikertelen
  hívásokat többször újrapróbálja, és szükség esetén újraindítja a modellt
  kiszolgáló alfolyamatot.
  \item \textbf{Dinamikus címkezelés.} A modellkiszolgáló címe futás közben
  frissíthető (\texttt{update\_model\_url}, illetve a \texttt{/model-url}
  végpont), a backend címe pedig a kliensoldali beállításokban módosítható, így
  a rendszer alkalmazkodik a változó nyilvános címekhez.
\end{itemize}

% =============================================================================
\section{Biztonsági megfontolások}
% =============================================================================

A rendszer a titkos kulcsokat (a Gemini API-kulcsot és a Firebase hozzáférési
adatokat) környezeti változókból, futásidőben tölti be, és \textbf{soha nem
tárolja a forráskódban}; e fájlokat a verziókezelés is kizárja. A backend minden
védett végponton ellenőrzi a Firebase azonosítótokent, így a klasszifikációs és
párbeszéd-funkciók kizárólag hitelesített felhasználók számára érhetők el. A
felhasználói adatok -- a feltöltött képek és a munkamenetek -- felhasználónként
elkülönítve tárolódnak, így egy felhasználó nem férhet hozzá más adataihoz.

Hangsúlyozandó továbbá, hogy a rendszer kialakítása tudatosan
\textbf{tájékoztató jellegű}: a klasszifikációt egy célzottan betanított,
kalibrált képi modell végzi, míg a nagy nyelvi modell szerepe a magyarázatra és a
párbeszédre korlátozódik. A rendszer nem helyettesíti a szakorvosi diagnózist,
hanem szűrőeszközként segíti a felhasználót a dermatológiai konzultáció
szükségességének megítélésében.
